% from intro-fr.tex
\oldpage[I]{9}
  déduit de $X$ par extension des scalaires au complété de $A$, à un résultat analogue pour une extension finie assez simple (par exemple non ramifiée) de $A$.
\end{enumerate}

Cette esquisse montre l'importance de l'étude systématique des schémas définis sur un anneau artinien $A$.
Le point de vue de Serre dans sa formulation de la théorie du corps de classes local, et des travaux récents de Greenberg, semblent suggérer qu'une telle étude pourrait être entreprise en attachant fonctoriellement à un tel schéma $X$ un schéma $X'$ sur le corps résiduel $k$ de $A$ (supposé parfait) de dimension égale (dans les cas favorables) à $n\dim X$, où $n$ est la longueur de $A$.

Quant à l'influence de A.~Weil, qu'il nous suffise de dire que c'est la nécessité de développer l'outillage nécessaire pour formuler avec toute la généralité voulue la définition de la <<~cohomologie de Weil~>> et pour aborder la démonstration\footnote{Pour éviter tout malentendu, précisons que cette tâche vient à peine d'être entreprise au moment où cette Introduction est écrite, et n'a donc pas encore abouti à la démonstration des conjectures de Weil.} de toutes les propriétés formelles nécessaires pour établir ses célèbres conjectures en Géométrie diophantienne~\cite{I-19}, qui a été une des principales motivations de la rédaction du présent Traité, au même titre que le désir de trouver le cadre naturel des notions et méthodes usuelles en Géométrie algébrique, et de donner aux auteurs l'occasion de comprendre lesdites notions et techniques.

\bigskip
\begin{center}*\quad*\quad*\end{center}
\bigskip

Pour terminer, nous croyons utile de prévenir les lecteurs que, tout comme les auteurs eux-mêmes, ils auront sans doute quelque difficulté avant de s'accoutumer au langage des schémas, et de se convaincre que les constructions habituelles que suggère l'intuition géométrique peuvent se transcrire, essentiellement d'une seule façon raisonnable, dans ce langage.
Comme dans beaucoup de parties de la Mathématique moderne, l'intuition première s'éloigne de plus en plus, en apparence, du langage propre à l'exprimer avec toute la précision et la généralité voulues.
En l'occurrence, la difficulté psychologique tient à la nécessité de transporter aux objets d'une catégorie déjà assez différente de la catégorie des ensembles (à savoir la catégorie des préschémas, ou la catégorie des préschémas sur un préschéma donné) des notions familières pour les ensembles~: produits cartésiens, lois de groupe, d'anneau, de module, fibrés, fibrés principaux homogènes, etc.
Il sera sans doute difficile au mathématicien, dans l'avenir, de se dérober à ce nouvel effort d'abstraction, peut-être assez minime, somme toute, en comparaison de celui fourni par nos pères, se familiarisant avec la Théorie des Ensembles.

\bigskip
\begin{center}*\quad*\quad*\end{center}
\bigskip

Les références seront données suivant le système décimal; par exemple, dans III, 4.9.3, le chiffre III indique le chapitre, le chiffre 4 le paragraphe, le chiffre 9 la section du paragraphe.
A l'intérieur du même chapitre, on supprimera la mention du chapitre.
